\documentclass[11pt]{article}
\usepackage[margin=1in]{geometry}
\usepackage{amsmath,amssymb}
\usepackage[T1]{fontenc}
\usepackage[utf8]{inputenc}
\usepackage{hyperref}

\title{Planar Bilipschitz Extension from Separated Nets: Literature Status}
\author{fa\_banach\_001, agent\_lane\_01}
\date{2026-07-18}

\begin{document}
\maketitle

\section*{Status}

\textbf{Literature already answered, two-dimensional subproblem.}
This note records that the planar case of the bilipschitz extension problem
for separated nets, flagged through Dymond--Kaluza arXiv:2507.22007, is
answered by the separate companion paper
\emph{Planar bilipschitz extension from separated nets}, Journal of the London
Mathematical Society 113 (2026), e70540.

\section*{Source Question}

Dymond and Kaluza, \emph{Extending Bilipschitz Mappings between Separated
Nets}, arXiv:2507.22007, recall the Alestalo--Trotsenko--Vaisala question:
for $d\geq 2$, does every bilipschitz mapping from $\mathbb Z^d$ to
$\mathbb R^d$, and more generally from a separated net of $\mathbb R^d$ to
$\mathbb R^d$, admit a bilipschitz extension to all of $\mathbb R^d$?
In the source, Theorem 1.1 reduces the arbitrary separated-net problem to the
integer-lattice problem in every finite dimension.

\section*{Answering Paper}

The companion article Dymond--Kaluza, \emph{Planar bilipschitz extension from
separated nets}, JLMS 113 (2026), e70540, DOI: 10.1112/jlms.70540, states as
Theorem 1.1 that every $L$-bilipschitz mapping
$f:\mathbb Z^2\to\mathbb R^2$ extends to a $p(L)$-bilipschitz mapping
$F:\mathbb R^2\to\mathbb R^2$ for a polynomial $p$. Combining this with the
reduction theorem from arXiv:2507.22007 gives Corollary 1.2: for every
separated net $X\subseteq \mathbb R^2$ and every $L$-bilipschitz
$f:X\to\mathbb R^2$, there is a bilipschitz extension
$F:\mathbb R^2\to\mathbb R^2$, with a bound depending polynomially on $L$ and
on the separation/net constants of $X$.

The supporting paper explicitly states that this answers Navas' 2015
Oberwolfach problem and gives the positive two-dimensional solution of the
Alestalo--Trotsenko--Vaisala problem.

\section*{Scope}

This packet does not claim any new proof. It records an already-known
literature answer for $d=2$. The higher-dimensional problem for $d\geq 3$,
including the hyperplane Conjectures 1 and 2 in arXiv:2507.22007, remains open
from the source's perspective.

\section*{Evidence}

The source arXiv PDF is included as \texttt{source\_paper.pdf}. The supporting
JLMS DOI page was checked on 2026-07-18; direct PDF download was blocked by a
Cloudflare challenge, saved as
\texttt{supporting\_paper\_jlms\_2026\_download\_blocked.html}. Web metadata
from Wiley and ISTA Research Explorer both record the 2026 JLMS publication and
the explicit planar theorem.

\begin{thebibliography}{9}
\bibitem{DK-extending}
M. Dymond and V. Kaluza,
\emph{Extending Bilipschitz Mappings between Separated Nets},
arXiv:2507.22007; Annales Fennici Mathematici 51 (2026), 237--260.

\bibitem{DK-planar}
M. Dymond and V. Kaluza,
\emph{Planar bilipschitz extension from separated nets},
Journal of the London Mathematical Society 113 (2026), e70540.
DOI: 10.1112/jlms.70540.
\end{thebibliography}

\end{document}
